% ---------------------------------------------------------------------------
% Error analysis (Section 6, first part) - drafted for M3 (Phase 8).
% Numbers come from notebooks/06_results_error_analysis.ipynb (Part 4) and
% notebook 07 (order stress test); figures from results/figures/.
% Followed by limitations_responsible_ai.tex. Citation keys match report/references.bib.
% ---------------------------------------------------------------------------
\section{Error Analysis and Limitations}
\label{sec:errors}

\subsection{Error analysis}
The analysis uses the test split and the final seed-42 runs, unless stated otherwise. Confusion rates for A3--A5 are
pooled over their three seeds.

\paragraph{The order-aware models use the order of the sounds}
On the test split, the anagram pair \textit{sita}/\textit{tisa} is the order-free A1's most-confused pair: 8.6\% of
the pair's clips are heard as the other word. For A2--A5 the rate falls to 1.0--1.9\%. Beyond that pair, spelling
similarity does not predict confusion. Across all 66 word pairs, the rank correlation between confusion rate and edit
distance lies between $-0.14$ and $+0.04$ for every model (all $p>0.2$). The closest pair, \textit{nane}/\textit{nne}
(edit distance 1), is confused in 4.7\% of its clips by A1, but in at most 0.3\% by the order-aware models.

The temporal-order stress test (Fig.~\ref{fig:stress}) checks the use of order directly. Each validation clip is
played unchanged, reversed, or cut into 100~ms chunks in random order. Every version keeps the same overall spectral
content, although reversal also reverses the dynamics within each sound, such as a stop's burst.
\begin{itemize}
  \item \emph{Reversal:} accuracy falls from 95--98\% to 21--39\% for A2--A5. On reversed audio, \textit{tisa} is
        mostly heard as \textit{sita} and vice versa: in 83\% of the pair's clips for the HMM and 62\% for A5.
  \item \emph{Control:} a logistic regression on static MFCC statistics alone is exactly unchanged by reversal
        (57.5\%).
  \item \emph{A1:} its delta features encode about 90~ms of local change, so it loses 28 points. It is order-free
        only beyond that scale.
  \item \emph{Shuffling:} A5 keeps 88.9\% accuracy when the chunks are shuffled, while A2--A4 fall to 32--61\%. So
        the pretrained model relies on the order \emph{within} chunks of about 100~ms (the transitions between
        sounds), and hardly on how those pieces are arranged. For the MFCC-based models the comparison is less
        clean, because the cuts also distort their delta features.
\end{itemize}

\begin{figure*}[t]
  \centering
  \includegraphics[width=0.78\linewidth]{figures/order_stress_test.png}
  \caption{Validation accuracy when each clip's word is played unchanged, reversed, or as shuffled 100~ms chunks.
  The control is A1 refitted on static MFCC statistics only.}
  \label{fig:stress}
\end{figure*}

\paragraph{The models fail on the same clips}
493 of the 630 test clips (78\%) are recognised by all five models, and 102 defeat only A1. Errors that belong to a
single order-aware architecture are rare: A5 makes none that A2--A4 all avoid, A4 one, A2 seven and A3 eight. The
same 10 test clips defeat A2, A3, A4 and A5 alike. They make up 10 of A5's 12 errors, and 10 of the 19--24 errors
of A2--A4.

Pooling those clips with the 11 validation clips that defeat A2--A5, we compared their recordings with all 1,260
validation and test clips (Table~\ref{tab:hard}).
\begin{itemize}
  \item \emph{Recording problems:} the hard clips are far more often very quiet, contain several separate
        utterances, or have speech at the edge of the default 1.5~s window, which suggests the crop chose the wrong
        stretch. Two-thirds of them are very quiet, noisy or cropped at the edge, against 18\% of all clips.
  \item \emph{Spectrograms:} in about three of the ten hard test clips, the window shows no clear word at all, only
        clicks or near-silence.
  \item \emph{Label-error candidates:} in three clips, A3--A5 agree on the same wrong word with at least 80\%
        confidence. One is a test clip labelled \textit{sita} that all four order-aware models hear as
        \textit{tisa}, with 99--100\% confidence.
\end{itemize}
Agreement among our models is not independent evidence of a label error, because they were all trained on these
labels. We therefore checked the labels in two independent ways.
\begin{itemize}
  \item \emph{Confident learning}~\cite{northcutt2021confident}, applied to the out-of-sample probabilities of A3--A5,
        flags only 1 of the 1,260 validation and test clips: the \textit{sita} clip above.
  \item \emph{An external recogniser:} the multilingual MMS model with its Swahili adapter~\cite{pratap2024scaling},
        which never saw this dataset, transcribed the 21 hard clips. On 60 easy test clips, it finds the labelled
        word in 82\% and never names a wrong one, so its positive findings can be trusted.
\end{itemize}
The recogniser explains six of the hard clips.
\begin{itemize}
  \item \emph{One likely label error:} it hears \textit{tisa} in the clip labelled \textit{sita}, in agreement with
        confident learning and our models.
  \item \emph{Three said in English:} two \textit{hapana} clips contain ``no'', and one \textit{tisa} clip contains
        ``nine''. Our models hear ``nine'' as \textit{nane}.
  \item \emph{Two crop failures:} the labelled word is in the recording, but outside the 1.5~s window.
\end{itemize}
The remaining transcripts contain other speech, such as a different Swahili phrase or English in the background, or
nothing. Because the recogniser also misses the word in 18\% of easy clips, these transcripts are weak evidence on
their own. Six of the 21 hard clips are therefore clear collection problems. For most of the rest, the recording
flags in Table~\ref{tab:hard} point to difficult audio rather than to a weakness of one architecture. Excluding the
two test clips that do not hold the labelled Swahili word raises every model's test accuracy by 0.3 points and leaves
the ranking unchanged. We did not relabel anything.

\begin{table}[t]
  \centering
  \caption{Recording flags among the 21 clips that every order-aware model misclassifies, against all 1,260
  validation and test clips. The window-edge flag uses the default 1.5~s window.}
  \label{tab:hard}
  \small
  \begin{tabular}{@{}lrr@{}}
    \toprule
    Flag & Hard clips & All clips \\
    \midrule
    Very quiet ($<-45$~dBFS) & 38\% & 4\% \\
    Several utterances & 52\% & 25\% \\
    Speech at the window edge & 43\% & 16\% \\
    Low SNR ($<15$~dB) & 14\% & 3\% \\
    Very short ($<0.3$~s of speech) & 10\% & 1\% \\
    Clipped & 0\% & 4\% \\
    \bottomrule
  \end{tabular}
\end{table}

\paragraph{Noise is the main remaining risk}
Error rates rise sharply in noisy recordings (Fig.~\ref{fig:strata}).
\begin{itemize}
  \item \emph{SNR:} below 30~dB (135 test clips), A2--A5 misclassify 8--13\% of clips, against at most 1.5\% above
        30~dB. A5 is the most robust (8.1\%). It is also the only model trained with noise injection, but this
        study cannot separate that from its pretraining.
  \item \emph{Loudness:} quiet recordings are harder for A2--A4 (4.7--8.3\% errors) than for A5 (2.4\%).
  \item \emph{Duration:} words estimated to last over 0.7~s show higher error rates. Part of this is noise
        inflating the estimate: duration and SNR are correlated ($\rho=-0.42$), and among the 50 long clips with good
        SNR, A5 makes one error. A2--A4 still make two or three there (4--6\%), above their overall rates.
\end{itemize}

\begin{figure*}[t]
  \centering
  \includegraphics[width=0.82\linewidth]{figures/p8_error_strata.png}
  \caption{Test error rate of the order-aware models by estimated word duration, SNR and loudness, with 95\%
  Wilson intervals.}
  \label{fig:strata}
\end{figure*}

\paragraph{What the models attend to}
\textit{nne} and \textit{nane} differ in the first half of the word, where \textit{nane} has the vowel /a/.
\begin{itemize}
  \item \emph{A3:} its attention peaks about a quarter of the way into both words. The profiles of the two words
        are indistinguishable ($p=0.79$). A bidirectional LSTM summarises the whole word in every state, so the
        attention shows where the model reads its summary, not which sounds decide.
  \item \emph{A4:} its class-activation evidence is spread over the whole word. For \textit{nane} it shifts slightly
        towards the start (56\% against 52\% in the first half, $p=0.014$), as expected if the vowel matters.
  \item \emph{Early recognition:} a forward-only A3 that sees only the beginning of the word reaches 24\% accuracy
        100~ms after the onset, 80\% at 500~ms and 88\% at 800~ms, against 93.7\% for the whole window (validation).
        The evidence separating the confusable words arrives with their second syllable.
\end{itemize}
Perturbation tests such as the stress test give firmer evidence about the use of order than these maps do.
